% c0_tensions.tex — §4

\section{\czero Tensions: \omh from All Probes}
\label{sec:c0_tensions}

\claim{c0_tensions_sign}{BAO alone prefers lower \omh ($+2.2\sigma$). All three SN datasets are mild or null: Union3 $-1.7\sigma$, Pantheon+ $-1.1\sigma$, DES-Dovekie $-0.8\sigma$ (all higher \omh or consistent with zero).}
\depends{c0_tensions_sign}{c0_universal}
\claim{bao_constrains_omh2}{BAO constrains \omh more tightly than the CMB ($\sigma_{c_0} = 1.56 > 1$). SN datasets do not ($\sigma_{c_0} \approx 0.28$--$0.33 < 1$).}
\depends{bao_constrains_omh2}{c0_universal}

Since all probes measure the same \czero direction (Section~\ref{sec:universal_c0}), we can compare their \czero measurements directly. The tension formula (Eq.~\ref{eq:tension}) gives four independent measurements of the same quantity, although we need to remember that the \czero normalization differs between probes. A unit shift in BAO \czero and a unit shift in SN \czero both correspond to $1\sigma$ of their respective measurement errors, but these translate to different shifts in \omh (see Section~\ref{sec:money_plot}).

Table~\ref{tab:c0_tensions} summarizes the results. A few points of interpretation are needed before examining the numbers:

\begin{itemize}
\item \textbf{Sign convention:} Positive tension means the data projection $a_0$ exceeds the chain mean $\langle c_0 \rangle$. Since $\beta_{\omh} = -0.90$, a positive \czero shift corresponds to \emph{lower} \omh (larger distances).

\item \textbf{What $\sigma_{c_0}$ means:} This is the standard deviation of \czero across the CMB chain posterior---it measures how much the CMB model prediction varies in the \czero direction, in units of measurement error. It is \emph{not} the measurement error, which is always 1 in whitened space by construction. The tension denominator $\sqrt{1 + \sigma_{c_0}^2}$ combines both.

\item \textbf{$\sigma_{c_0}$ and constraining power:} BAO has $\sigma_{c_0} = 1.56 > 1$. This means BAO constrains \omh more tightly than the CMB: DESI data carries genuinely new information about \omh beyond what the CMB already provides. The three SN datasets have $\sigma_{c_0} \approx 0.28$--$0.33 < 1$: the CMB already provides a tighter constraint than these SN measurements, so their \czero tensions are measured against a prior that is already narrower than the data.
\end{itemize}

\begin{deluxetable*}{lcccc}
\tablecaption{\czero Tension per Probe \label{tab:c0_tensions}}
\tablecomments{Columns: $\langle c_0 \rangle$ is the mean of the CMB chain's \czero distribution (the model prediction in the $V_0$ direction); $\sigma_{c_0}$ its standard deviation across the chain (model uncertainty, in units of the measurement error); $a_0$ is the data amplitude, the projection of the measured low-$z$ data onto $V_0$ (measurement error 1 by construction); and the tension is $(a_0 - \langle c_0 \rangle)/\sqrt{1+\sigma_{c_0}^2}$ (Eq.~\ref{eq:tension}). Values in parentheses use the Planck 2018 \LCDM chain as the CMB reference instead of ACT DR6.}
\tablehead{
\colhead{Probe} & \colhead{$\langle c_0 \rangle$} & \colhead{$\sigma_{c_0}$} & \colhead{$a_0$} & \colhead{Tension}
}
\startdata
BAO (DESI) & $+0.59$ ($-0.31$) & 1.56 (1.85) & $+4.69$ & $+2.2\sigma$ ($+2.4\sigma$) \\
Union3 & $+0.13$ ($-0.06$) & 0.28 (0.34) & $-1.59$ & $-1.7\sigma$ ($-1.4\sigma$) \\
Pantheon+ & $+0.14$ ($-0.06$) & 0.30 (0.36) & $-0.97$ & $-1.1\sigma$ ($-0.9\sigma$) \\
DES-Dovekie & $+0.16$ ($-0.07$) & 0.33 (0.40) & $-0.74$ & $-0.8\sigma$ ($-0.6\sigma$) \\
\enddata
\end{deluxetable*}

\evidence{c0_tensions_sign}{ev_c0_table}{Four-probe tension table: BAO positive, all SN mild/negative.}

BAO prefers lower \omh than the CMB (positive \czero tension at $+2.2\sigma$), while the three SN datasets prefer higher \omh --- Union3 at $-1.7\sigma$, Pantheon+ at $-1.1\sigma$, and DES-Dovekie at $-0.8\sigma$. None of the SN tensions is individually significant, and because $\sigma_{c_0}^\mathrm{SN} < 1$ the SN datasets do not have the constraining power on \czero to confirm or refute the BAO result; the opposite-sign pattern is suggestive but should not be interpreted as a measured disagreement.

An important point to remember when interpreting Figure~\ref{fig:c0_gaussians}: the \czero normalization differs between probes. A unit shift in BAO \czero and a unit shift in SN \czero both correspond to $1\sigma$ of their respective measurement errors, but these translate to different shifts in \omh (see Section~\ref{sec:money_plot}). The figure is designed to visualize \emph{tensions}, not to compare \omh values directly. 

\begin{figure}[t]
\centering
\includegraphics[width=\columnwidth]{figures/fig_c0_gaussians.pdf}
\caption{Per-probe \czero distributions. \emph{Shaded regions}: CMB chain predictions (model uncertainty in each probe's \czero). \emph{Dashed curves}: measurement likelihoods centered at the data \czero value ($\sigma = 1$ in whitened space by construction). The BAO chain distribution (shaded, wide) has $\sigma_{c_0} = 1.56$, wider than the measurement likelihood, reflecting BAO's ability to constrain \omh beyond the CMB prior. The three SN chain distributions (shaded, narrow, $\sigma_{c_0} \approx 0.28$--$0.33$) are much narrower than their measurement likelihoods, indicating SN data adds little constraining power on \czero.}
\label{fig:c0_gaussians}
\end{figure}

\subsection{The \omh Money Plot}
\label{sec:money_plot}

Since \czero is a near-perfect proxy for \omh (Eq.~\ref{eq:c0_formula}), we can convert each probe's \czero measurement to an implied \omh value. This is the value of \omh implied by the data if one slides along the CMB posterior's degeneracy direction in $(\omh, \obh, \thetastar)$. The conversion inverts the $\beta$-vector formula, setting \obh and \thetastar to their chain mean values:
\begin{equation}
\omh_\mathrm{probe} = \langle \omh \rangle_\mathrm{chain} + (a_0 - \langle c_0 \rangle) \times \frac{\sigma_{\omh}}{\beta_{\omh} \cdot \sigma_{c_0}}
\label{eq:omh_conversion}
\end{equation}
where $\beta_{\omh}$ is the $\beta$-vector coefficient for \omh, and $\sigma_{\omh}$, $\sigma_{c_0}$ are the chain standard deviations. The measurement error on \omh from the $\sigma = 1$ whitened-space uncertainty is $\sigma_{\omh}^\mathrm{meas} = |\sigma_{\omh} / (\beta_{\omh} \cdot \sigma_{c_0})|$. An additional marginalization uncertainty arises from the CMB posterior widths of \obh and \thetastar; this adds $\sigma_\mathrm{marg} \approx 0.0003$ in quadrature, a modest correction ($\sim 7\%$ for BAO, negligible for SN).

Figure~\ref{fig:omegamh2_money} shows the result. BAO provides the tightest \omh constraint of the four probes --- tighter than the CMB itself --- while the three SN measurements have total errors $\sim 5\times$ larger.

\begin{figure}[t]
\centering
\includegraphics[width=\columnwidth]{figures/fig_omegamh2_money.pdf}
\caption{Effective \omh measurements from all four distance probes, derived from \czero via the $\beta$-vector inversion (Eq.~\ref{eq:omh_conversion}). Error bars show total uncertainty (measurement $+$ marginalization over \obh and \thetastar, added in quadrature). The gray band shows the ACT \LCDM posterior width. BAO provides a tighter constraint than the CMB.}
\label{fig:omegamh2_money}
\end{figure}

\subsection{Physical Interpretation of the Mode Shape}

The \czero mode measures \omh, and the physical effect of varying \omh at fixed \thetastar is a characteristic tilt of the Hubble rate $H(z)$ (Figure~\ref{fig:hz_omh2}). Increasing \omh raises $H(z)$ at high redshift (where $H \propto \sqrt{\omh}$ in the matter era), but the \thetastar constraint forces $H_0$ to \emph{decrease} to compensate (keeping $r_\star / D_A(z_\mathrm{rec})$ fixed). The result is a crossover near $z \sim 1$--$2$: higher \omh means faster expansion at $z > 2$ but slower at $z < 1$. This crossover falls squarely in the DESI BAO redshift range, explaining why BAO is sensitive to \omh.

For distances, lower \omh (positive \czero) produces smaller $D_V/\rs$ ratios at all BAO redshifts (Figure~\ref{fig:c0_families}, left panel): the slower high-$z$ expansion reduces distances, while the higher $H_0$ reduces $\rs$-normalized quantities further. For SN, the same physics manifests as a tilt in $\mu(z)$: positive \czero shifts distance moduli downward (brighter supernovae) at low $z$ and upward at high $z$. Note that when making these plots we just varied \omh while keeping \obh and \thetastar fixed for simplicity. Varying \obh and \thetastar along the degeneracy does not change things qualitatively.

\begin{figure*}[t]
\centering
\includegraphics[width=\textwidth]{figures/fig_hz_omh2_variation.pdf}
\caption{Hubble rate ratio $H(z)/H_\mathrm{ref}(z)$ for varying \omh at fixed \thetastar (flat \LCDM). Higher \omh (blue, dashed) raises $H$ at high $z$ but lowers $H_0$; the crossover near $z \sim 2$ falls within the DESI BAO range (shaded). The red curve shows the BAO-implied \omh from the \czero measurement ($\omh = 0.139$, $H_0 = 70.6$).}
\label{fig:hz_omh2}
\end{figure*}

\subsection{Goodness of Fit and Robustness}
\label{sec:robustness}

Beyond the \czero tension, a natural question is whether the remaining SVD modes are well-behaved. We compute a dimension-limited $\chi^2$:
\begin{equation}
\chi^2(K) = \sum_{\alpha=0}^{K-1} \frac{(a_\alpha - \langle c_\alpha \rangle)^2}{1 + \sigma_{c_\alpha}^2}
\label{eq:chi2_K}
\end{equation}
For BAO, $\chi^2_{13} = 15.7$ ($p = 0.27$); excluding \czero, $\chi^2_{12} = 10.8$ ($p = 0.55$). The remaining 12 modes are well-behaved. All three SN datasets also show good fits: Union3 $\chi^2_{22} = 26.8$ ($p = 0.22$), Pantheon+ $\chi^2_{22} = 28.4$ ($p = 0.16$), DES-Dovekie $\chi^2_{19} = 20.0$ ($p = 0.39$). Excluding \czero leaves Union3 $\chi^2_{21} = 24.0$ ($p = 0.29$), Pantheon+ $\chi^2_{21} = 27.3$ ($p = 0.16$), DES-Dovekie $\chi^2_{18} = 19.3$ ($p = 0.37$). No dataset shows anomalous higher-mode structure. Of course this is a very crude test, as it could be that a significant discrepancy in one specific direction would get diluted in this statistic. Thus in the next sections we will perform a more detailed analysis of the higher-mode structure using specific physical models as guidance to identify promising directions based on our theoretical expectations. 

We also verify that the tensions are robust to the redshift range: restricting all datasets to $z \leq 0.9$ changes tensions modestly (BAO: $+2.2\sigma \to +1.9\sigma$; Union3: $-1.7\sigma \to -2.0\sigma$; Pantheon+: $-1.1\sigma \to -1.8\sigma$; DES-Dovekie: $-0.8\sigma \to -0.6\sigma$). The high-redshift bins do not drive the tensions.

\subsection{Summary Within \LCDM}
\label{sec:central_finding_lcdm}

Within \LCDM, there is a tension between DESI and the CMB which lives in the direction we are calling the \omh direction---the worst-constrained CMB combination parameter for low-$z$ expansion history. No other direction is measurable: $\sigma_{c_1} = 0.06$ for BAO, $\approx 0.001$ for SN (Table~\ref{tab:sigma_c}). This fact is not particularly surprising. The BAO \czero tension is mild but real: BAO and the CMB prefer different \omh values at $+2.2\sigma$. The three SN datasets do not have enough constraining power in $V_0$ ($\sigma_{c_0}^\mathrm{SN} < 1$) to either confirm or refute the BAO signal at its precision; their tensions are mild or null but cannot independently corroborate a true cosmological shift.
In the next section, we extend the analysis to \wowa dark energy and ask whether this tension opens genuinely new directions, or whether the DESI \wowa preference reduces primarily to the \czero tension once the data are decomposed in the \LCDM SVD basis.

\dataref{ev_c0_table}{calculation/scripts/cross\_dataset\_analysis.py}
\dataref{ev_c0_table}{calculation/scripts/diagnose\_union3.py}
